\documentclass[11pt,letterpaper]{article}
\usepackage[T1]{fontenc}
\usepackage[utf8]{inputenc}
\usepackage{lmodern,microtype}
\usepackage[margin=1in]{geometry}
\usepackage{amsmath,amssymb,amsthm,mathtools}
\usepackage{booktabs,tabularx,array,longtable}
\usepackage{xurl}
\usepackage[hidelinks]{hyperref}
\usepackage{enumitem}
\hypersetup{pdftitle={Signed Mellin Curvature and Theta Coefficient Positivity},pdfauthor={Drafted for Edward Baker}}
\setlist{itemsep=3pt,topsep=5pt}
\setlength{\emergencystretch}{2em}
\setlength{\parskip}{3pt}
\usepackage{tocloft}
\setlength{\cftbeforesecskip}{2pt}
\setcounter{tocdepth}{1}
\numberwithin{equation}{section}
\newtheorem{proposition}{Proposition}[section]
\newtheorem{lemma}[proposition]{Lemma}
\newtheorem{hypothesis}[proposition]{Research hypothesis}
\theoremstyle{remark}
\newtheorem{remark}[proposition]{Remark}
\DeclareMathOperator{\Var}{Var}
\DeclareMathOperator{\Cov}{Cov}
\DeclareMathOperator{\Tr}{Tr}
\DeclareMathOperator{\diag}{diag}
\newcommand{\E}{\mathbb E}
\newcommand{\R}{\mathbb R}
\newcommand{\N}{\mathbb N}
\newcommand{\dd}{\,d}
\newcommand{\evid}[1]{\textsf{[#1]}}
\newcommand{\ket}[1]{\lvert #1\rangle}
\newcommand{\bra}[1]{\langle #1\rvert}
\newcommand{\ip}[2]{\langle #1,#2\rangle}
\title{Signed Mellin Curvature and\newline Theta Coefficient Positivity\\[5pt]
\large Mathematical Background and Physical Research Directions}
\author{Drafted for Edward Baker}
\date{2 November 2026}

\begin{document}
\maketitle
\begin{abstract}
The Mellin-curvature program seeks a quantitative way to raise the order
of positivity of the Taylor-coefficient determinants associated with the
Riemann xi function. Its starting point is an exact integral in which a
positive theta carrier is multiplied by a sign-changing observable.
At determinant order three, an internally assessed signed-averaging
argument proves positivity for every nonnegative real Mellin tilt. The
remaining question is an upper bound on the logarithmic curvature of the
complete signed integral. We derive the factorial comparison that turns
such a bound into an order-four inequality, explain the Gaussian
normalization, and separate this sufficient conjecture from the exact
next-order sign condition and the all-order Riemann-hypothesis criterion.
We then describe eight complementary investigations: thermal traces,
fermionic amplitudes, connected response, biorthogonal ensembles,
Schur and Toda identities, supersymmetric protection, four-dimensional
gauge-theory anomalies, and open quantum dynamics. Exact physical
realizations, local asymptotics, and scoped obstructions are distinguished
from the still-missing signed estimate. The purpose is to provide shared
mathematical background and a common specification of what a successful
subinvestigation must establish.
\end{abstract}

\noindent\textbf{Status.} Expository research draft, incorporating the records
available on 2 November 2026. No new domain of the Gaussian Mellin
comparison is established here. Internal LLM assessments cited below are
not independent specialist review. Source identifiers such as \evid{M1}
are resolved in Appendix~\ref{sec:sources}.

\tableofcontents
\clearpage

\section{The question and its place in the Jensen program}
\label{sec:question}

The program begins with a simple tension. The canonical theta source is
positive, and its even moments are positive. Yet the coefficient
determinants relevant to real zeros contain alternating sums. Their
natural integral representation has a genuinely negative region. The
aim is to exploit the structure of that signed integral rather than
replace it by a positive object with a different meaning.

Two indices must be kept distinct. Increasing the coefficient index
changes which part of the source is emphasized by a moment weight.
Increasing the determinant order introduces more source points and
higher derivatives. An estimate that handles the first operation need
not handle the second. Mellin curvature is proposed as the extra
quantitative information that could connect them.

\subsection{The complete source and its normalization}

For $\beta>0$, define
\begin{equation}
 \vartheta(\beta)=\sum_{m\in\mathbb Z}e^{-\pi\beta m^2},
 \qquad p(u)=\tfrac12e^u\bigl(\vartheta(e^{4u})-1\bigr).
 \label{eq:primitive}
\end{equation}
The source used throughout is
\begin{align}
 \Phi(u)&=\tfrac18\bigl(p''(u)-p(u)\bigr) \notag\\
 &=\sum_{m\ge1}\bigl(2\pi^2m^4e^{9u}-3\pi m^2e^{5u}\bigr)
                 e^{-\pi m^2e^{4u}}. \label{eq:source}
\end{align}
Every summand is positive for $u\ge0$. The complete function is even,
by theta inversion, and its fixed derivatives decay double
exponentially as $u\to+\infty$. Evenness is a property of the complete
source; it must not be attributed to a finite truncation of this series.

Set
\begin{equation}
 b_n=\int_0^\infty u^{2n}\Phi(u)\dd u,\qquad
 c_n=\frac{b_n}{(2n)!},\qquad \gamma_n=n!c_n,
 \qquad F(w)=\sum_{n\ge0}c_nw^n. \label{eq:coeff}
\end{equation}
These are respectively raw moments, ordinary coefficients, and
exponential coefficients. They are not interchangeable. With
$\xi(z)=\tfrac12z(z-1)\pi^{-z/2}\Gamma(z/2)\zeta(z)$, the convention is
\begin{equation}
 F(z^2)=\int_0^\infty\Phi(u)\cosh(zu)\dd u
       =\tfrac18\xi\bigl((1+z)/2\bigr). \label{eq:xi}
\end{equation}
The factor $1/8$ and the rescaling of $z$ matter when comparing sources.
The integral defines an entire function: on compact $z$ sets its growth
factor is dominated by the source decay. The classical theta Mellin
representation underlies \eqref{eq:xi}; the repository's normalization is
retained in \evid{J1,M3,G1}.

\subsection{The full destination and the finite-order pilot}

The Jensen polynomials and consecutive Toeplitz determinants are
\begin{equation}
 J^{d,n}(X)=\sum_{j=0}^d\binom dj\gamma_{n+j}X^j,
 \qquad D_{r,k}=\det[c_{k+j-i}]_{i,j=0}^{r-1}, \label{eq:targets}
\end{equation}
where $c_j=0$ for negative integers and $D_{0,k}=1$.
Hyperbolicity means that every zero of a polynomial is real.
The complete classical destination is
\begin{equation}
 \mathrm{RH} \Longleftrightarrow\ F\text{ has only negative real zeros}
 \Longleftrightarrow\ (c_n)\in\mathrm{PF}_{\infty}
 \Longleftrightarrow\ J^{d,n}\text{ is hyperbolic for every }d,n.
 \label{eq:rh}
\end{equation}
Here $\mathrm{PF}_{\infty}$ means that every finite minor of the
infinite Toeplitz matrix $[c_{j-i}]_{i,j\ge0}$ is nonnegative.
The classical entire-function/P\'olya-frequency bridge is stated in
Katkova's introduction \cite{Katkova0505174v1}; the Jensen formulation
and its quantifiers are recalled in \cite{GriffinOnoRolenZagier2019Jensen}.
Rescaling $c_n$ to $A B^n c_n$, with fixed $A,B>0$, does not alter
these criteria.

Strict positivity of the consecutive minors $D_{r,k}$ for all $r,k$ is
a sufficient route to the full minor condition. The strict-minor lemma
quoted by Katkova requires a positive summable sequence and strict
positivity at each order up to the intended order. Here $c_n>0$ and
$\sum c_n=F(1)<\infty$. A bare nonnegative consecutive-minor statement
must not silently replace that strict hypothesis.

For Jensen polynomials, fixed degree with sufficiently large shift is
already covered by the eventual hyperbolicity theorem of
Griffin--Ono--Rolen--Zagier. Its threshold may depend on the degree.
It does not cover unbounded degree at a fixed small shift. In fact
\begin{equation}
 \frac{d^n}{dX^n}J^{d+n,0}(X)=\frac{(d+n)!}{d!}J^{d,n}(X),
\end{equation}
so all degrees at shift zero would suffice. The determinant anchor
$D_{r,0}=c_0^r$ is a different statement; it does not solve this
all-degree problem.

Known finite-order positivity already includes order four; for example,
Katkova records a $\mathrm{PF}_{44}$ result for the rescaled xi
coefficient function \cite[Theorem 1]{Katkova0505174v1}. The present
$r=3\to4$ task is therefore a \emph{method pilot}. Its value would be a
new quantitative mechanism and a credible route to arbitrary order,
not a newly discovered low-order sign range.

\section{From coefficient determinants to a signed average}
\label{sec:integral}

\subsection{An exact integral, including the small-index endpoints}

Let
\[
 \Omega_r=\{0<u_1<\cdots<u_r\},\qquad
 V(\mathbf y)=\prod_{a<b}(y_b-y_a),\qquad
 \Delta_r(\mathbf u)=\det[\Phi^{(2i)}(u_j)]_{i=0,\ldots,r-1;\ j=1,\ldots,r}.
\]
Define the factorial normalization separately from the generating
function $F$:
\begin{equation}
 F_{r,k}=\prod_{j=0}^{r-1}(2k+2j)!,\qquad
 M_{r,k}=F_{r,k}D_{r,k}. \label{eq:factorial}
\end{equation}
Then the basic representation is
\begin{equation}
 M_{r,k}=\int_{\Omega_r}\Bigl(\prod_j u_j^{2k}\Bigr)
                    V(\mathbf u^2)\Delta_r(\mathbf u)\dd\mathbf u.
 \label{eq:integral}
\end{equation}

To see where every factor comes from, integration by parts gives
\begin{equation}
 \int_0^\infty\frac{u^{2m}}{(2m)!}\Phi^{(2i)}(u)\dd u=c_{m-i}.
 \label{eq:ibp}
\end{equation}
For $m\ge i$, transfer $2i$ derivatives to the polynomial. For $m<i$,
transfer $2m$ derivatives first; the remaining integral is
$-\Phi^{(2i-2m-1)}(0)=0$. Thus the zero convention for negative
coefficient indices follows from \emph{all} odd source derivatives
vanishing at zero. The upper endpoints vanish by decay.
The determinant-composition identity applied to \eqref{eq:ibp}
produces a factor $1/r!$ on the full orthant. Both determinants change
sign under a node permutation, so their product is symmetric. Restricting
to $\Omega_r$ cancels that factor. The monomial determinant supplies
$V(\mathbf u^2)$ and the column factorials supply $F_{r,k}$.
This proves \eqref{eq:integral} with no extra chamber constant.

At $k=0$ the coefficient matrix is upper triangular:
\begin{equation}
 D_{r,0}=c_0^r. \label{eq:anchor}
\end{equation}
The first substantive neighboring-index problem is consequently $k=1$.
Large-$k$ localization does not remove the need to handle it.

\subsection{What is positive and what remains signed}

The accepted order-three analysis \evid{J2} establishes
$R'(u)>0$ for $u>0$, where $R=\Phi''/\Phi$.
Writing $x=R(u)$, put
\begin{equation}
 g_i(x)=\left.\frac{\Phi^{(2i)}}{\Phi}\right|_{u=R^{-1}(x)},
 \qquad
 h_r(\mathbf u)=\frac{\det[g_i(R(u_j))]_{i,j}}{V(R(\mathbf u))}.
 \label{eq:hr}
\end{equation}
The quotient has removable coincident-node limits. Since $g_0=1$ and
$g_1=x$, one has $h_2=1$, while $h_3$ is the second divided difference
of $g_2$. Define
\begin{align}
 w_{r,s}(\mathbf u)&=\Bigl(\prod_j u_j^{2s}\Phi(u_j)\Bigr)
                     V(\mathbf u^2)V(R(\mathbf u)),\notag\\
 Z_r(s)&=\int_{\Omega_r}w_{r,s},\qquad
 d\mu_{r,s}=Z_r(s)^{-1}w_{r,s}\dd\mathbf u,\notag\\
 A_r(s)&=\E_{r,s}h_r,\qquad
 M_r(s)=Z_r(s)A_r(s)=\int_{\Omega_r}h_rw_{r,s}. \label{eq:carrier}
\end{align}
The carrier is positive. The insertion $h_r$ need not be.
At integer $s=k$, this is precisely \eqref{eq:integral}.
For noninteger $s$ it is an interpolation of integrals; it is not a
Toeplitz determinant with noninteger coefficient indices.

For order three, the additional global derivative sign assessed in
\evid{J2} implies $g_2'''(x)>0$ in the interior. The simplex formula
for divided differences writes $h_3$ as the integral of $g_2''$ over a
triangle of area $1/2$. Increasing $g_2''$ therefore makes $h_3$
increasing in each node. This is compatible with the established
negative values of $g_2''$ near the origin and the resulting negative
region of $h_3$. The proof needs monotonicity, not pointwise positivity.

The carrier's Vandermonde factors satisfy the multiplicative lattice
inequality on the ordered chamber. Finite-grid positive association,
followed by an integrable exhaustion, then gives, for $s,\ell\ge0$,
\begin{equation}
 A_3(s+\ell)-A_3(s)
 =\frac{\Cov_{3,s}(h_3,q^\ell)}{\E_{3,s}q^\ell}\ge0,
 \qquad q=\prod_{j=1}^3u_j^2. \label{eq:association}
\end{equation}
The integer version has an internal certificate/assessment record;
the real-tilt extension is justified in \evid{M1} by the same
association and exhaustion argument. In particular,
\begin{equation}
 A_3(s)\ge A_3(0)=\frac{48c_0^3}{Z_3(0)}>0,
 \qquad M_3(s)>0\quad(s\ge0). \label{eq:realpositive}
\end{equation}
This separately justifies the logarithm used below. Integer determinant
positivity alone would not justify it between the integers.

All fixed-order integrals with finitely many logarithmic insertions are
absolutely convergent. Near zero, use the undivided expression
$h_rw_{r,s}=(\prod u_j^{2s})V(\mathbf u^2)\Delta_r$ and bounded source
derivatives; $u^{2s}|\log u|^m$ is locally integrable for $s\ge0$.
At infinity, double-exponential decay dominates all fixed derivative,
moment, and logarithmic factors. These observations also explain why
coincident-node quotients introduce no missing boundary term.

\section{The complete Mellin curvature}
\label{sec:curvature}

From now on, unless an order is displayed, set $r=3$ and abbreviate
$M=M_3$, $Z=Z_3$, $A=A_3$, $h=h_3$. The real variable $s$ is a
\emph{Mellin tilt}. It is not the zeta argument, thermal inverse
temperature, or a physical evolution time. Put
\begin{equation}
 T=\log q=2\sum_{j=1}^3\log u_j,\qquad m_s=\E_sT.
\end{equation}
The weight $q^s=e^{sT}$ makes $s$ a source coupled to the logarithmic
size of the three-node configuration.

\begin{proposition}[Signed connected response]
For the complete integral above,
\begin{equation}
 \mathcal C_3(s):=(\log M)''(s)
 =\Var_s(T)
 +\frac{\Cov_s\bigl(h,(T-m_s)^2\bigr)}{A(s)}
 -\left(\frac{\Cov_s(h,T)}{A(s)}\right)^2,
 \qquad s\ge0. \label{eq:curvature}
\end{equation}
At the endpoint, derivatives may equivalently be interpreted as
right derivatives.
\end{proposition}
\begin{proof}
Differentiation under the integral gives
$M'=\int Thw_s$ and $M''=\int T^2hw_s$.
Thus $(\log M)''=\E_s(hT^2)/A-(\E_s(hT)/A)^2$.
Subtract $m_s$ from $T$ in this expression, expand, and use
$\E_s[h(T-m_s)]=\Cov_s(h,T)$. The preceding integrability bounds
justify the differentiations locally uniformly in $s$.
\end{proof}

The first term in \eqref{eq:curvature} is an ordinary variance under a
positive measure. The other two terms record the insertion's response.
The complete expression is not generally a variance under a positive
probability law: $h\,d\mu_s/A$ is signed. In particular, no generic
nonnegativity of $\mathcal C_3$ is asserted.

Association controls $\Cov_s(h,T)$, since both functions are increasing.
It does not control $\Cov_s(h,(T-m_s)^2)$: a centered square is not
increasing over the entire range of $T$. Bounding the carrier variance
alone leaves this contribution unestimated. The helpful negative square
in \eqref{eq:curvature} is also part of the same expression. Separating
all signed contributions by absolute values can lose exactly the
cancellation needed for the desired upper bound.

An alternative form makes the issue especially transparent. If
$f_s=h w_s$, then
\begin{equation}
 M(s)^2\mathcal C_3(s)
 =\frac12\int_{\Omega_3^2}
 f_s(\mathbf u)f_s(\mathbf v)
 \bigl(T(\mathbf u)-T(\mathbf v)\bigr)^2\dd\mathbf u\dd\mathbf v.
 \label{eq:twocopycurv}
\end{equation}
Pairs drawn from opposite signs of the insertion contribute negatively.
This is an exact two-copy identity, not a positivity argument.

\section{The order-raising budget and the Gaussian hypothesis}
\label{sec:budget}

\subsection{The exact budget}

Desnanot--Jacobi applied to a Toeplitz block gives
\begin{equation}
 D_{r+1,k}D_{r-1,k}=D_{r,k}^2-D_{r,k-1}D_{r,k+1}.
 \label{eq:dj}
\end{equation}
For $k\ge1$, suppose that the displayed order-$r$ determinants and
$D_{r-1,k}$ are positive, and define
\begin{equation}
 Q_{r,k}=\frac{M_{r,k-1}M_{r,k+1}}{M_{r,k}^2},\qquad
 B_{r,k}=\frac{F_{r,k-1}F_{r,k+1}}{F_{r,k}^2}.
\end{equation}
The normalization telescopes:
\begin{align}
 B_{r,k}
 &=\prod_{j=0}^{r-1}
 \frac{(2k+2j+2)(2k+2j+1)}{(2k+2j)(2k+2j-1)}\notag\\
 &=\frac{(2k+2r)(2k+2r-1)}{2k(2k-1)}. \label{eq:B}
\end{align}
Consequently
\begin{equation}
 D_{r+1,k}=\frac{D_{r,k}^2}{D_{r-1,k}}
                  \left(1-\frac{Q_{r,k}}{B_{r,k}}\right),
 \qquad D_{r+1,k}>0\ \Longleftrightarrow\ Q_{r,k}<B_{r,k}.
 \label{eq:exactsign}
\end{equation}
The last inequality is exactly the next determinant's sign condition.
It is a reformulation, not an estimate acquired by the program.
For fixed $r$, $B_{r,k}=1+2r/k+O_r(k^{-2})$: the permitted excess
over one shrinks at large coefficient index.

\subsection{A stronger comparison with a quantitative payoff}

The selected sufficient comparison is
\begin{equation}
 Q_{r,k}\le G_{r,k}:=\frac{2k+2r-1}{2k-1}. \label{eq:G}
\end{equation}
It leaves $B_{r,k}/G_{r,k}=(k+r)/k$. Hence, whenever the lower-order
positivity assumptions hold,
\begin{equation}
 D_{r+1,k}\ge\frac{r}{k+r}\frac{D_{r,k}^2}{D_{r-1,k}}>0.
 \label{eq:payoff}
\end{equation}
At the pilot order $r=3$, $D_{2,k}>0$ follows independently from
$h_2=1$, and the neighboring $D_{3,k}$ are positive by
\eqref{eq:realpositive}. No order-four sign is needed to derive this
conditional implication.

\begin{hypothesis}[Pointwise Gaussian curvature envelope]\label{hyp:gaussian}
For the unchanged canonical theta source at physical heat time zero,
\begin{equation}
 \mathcal C_3(s)\le\kappa_3(s),\qquad
 \kappa_r(s):=\sum_{j=0}^{r-1}\sum_{\nu=0}^\infty
                    (s+j+\tfrac12+\nu)^{-2},\qquad s\ge0.
 \label{eq:hypothesis}
\end{equation}
\end{hypothesis}
This hypothesis remains unproved. Its role is a specific, falsifiable
sufficient estimate; neither its statement nor an exact physical
representation supplies evidence for its validity.

\begin{proposition}[The pilot implication]
Hypothesis~\ref{hyp:gaussian} implies, for every integer $k\ge1$,
\begin{equation}
 Q_{3,k}\le\frac{2k+5}{2k-1},\qquad
 D_{4,k}\ge\frac{3}{k+3}\frac{D_{3,k}^2}{D_{2,k}}>0.
 \label{eq:pilot}
\end{equation}
\end{proposition}
\begin{proof}
For a twice differentiable function $L$,
\[
 L(k-1)+L(k+1)-2L(k)
 =\int_{k-1}^{k+1}(1-|s-k|)L''(s)\dd s.
\]
Apply this to $L=\log M$. For $b=k+j+1/2+\nu>1$, the corresponding
summand of the integral of $\kappa_r$ is
$\log\bigl(b^2/(b^2-1)\bigr)$. The product over $\nu$ telescopes,
and then the product over $j$ telescopes, giving
\begin{equation}
 \int_{k-1}^{k+1}(1-|s-k|)\kappa_r(s)\dd s
 =\sum_{j=0}^{r-1}\log\frac{k+j+1/2}{k+j-1/2}
 =\log G_{r,k}. \label{eq:triangle}
\end{equation}
The series is positive and convergent, justifying integration termwise.
Exponentiation proves the first assertion; \eqref{eq:payoff} proves the
second.
\end{proof}

The logical hierarchy is important:
\[
 \text{pointwise curvature envelope}
 \ \Longrightarrow\ \text{integrated Gaussian comparison}
 \ \Longrightarrow\ \text{next-order positivity}.
\]
A failure of the first condition would not refute the second, and a
failure of the second would not refute order-four positivity. At $k=1$
the required integration interval is exactly $[0,2]$; neither negative
Mellin tilts nor negative-index moments enter.

\subsection{Why the comparator is Gaussian}

For the comparison source $\Phi_G(u)=e^{-au^2}$, $a>0$,
\[
 c_n^G=c_0^G\frac{t^n}{n!},\qquad
 c_0^G=\frac{\sqrt\pi}{2\sqrt a},\qquad t=\frac1{4a}.
\]
Factoring the reciprocal-factorial matrix into falling factorials gives
\begin{equation}
 D^G_{r,k}=(c_0^G)^r t^{rk}
                  \prod_{j=0}^{r-1}\frac{j!}{(k+j)!}. \label{eq:gaussD}
\end{equation}
The derivative ratios are monic polynomials of degree $i$ in $R$,
so $h_r^G=1$. Equivalently, changing variables to $u_j^2$ in the
positive Vandermonde-squared integral yields
\begin{equation}
 M_r^G(s)=C_{r,a}\,a^{-rs}
                 \prod_{j=0}^{r-1}\Gamma(s+j+\tfrac12),
 \qquad C_{r,a}>0. \label{eq:gaussM}
\end{equation}
Thus $(\log M_r^G)''=\kappa_r$ and $Q^G_{r,k}=G_{r,k}$.
This is an exact normalization check. Theta has a different tail and a
signed insertion; Gaussian equality does not prove its upper bound.

\subsection{Why association alone cannot supply the budget}

The exact relaxation in \evid{M1} uses two values $q=1,128$, base
weights $128,1$, and increasing observable values $h=1,2$. It has
\[
 Z(s)=128+128^s,\qquad M(s)=128+2\cdot128^s,\qquad A(s)=M(s)/Z(s).
\]
Association holds, $A$ is positive and increasing, and $Z$ is a positive
moment function. Nevertheless,
\begin{equation}
 \frac{M(0)M(2)}{M(1)^2}=\frac{16705}{576}>28=B_{3,1}>7=G_{3,1}.
\end{equation}
These assumptions do not even ensure the weaker factorial budget.
The example is a relaxation of available information: it is not the
canonical source, not a derivative-determinant construction, and not
a Toeplitz counterexample. It identifies the need for additional
theta-specific structure.

\section{Eight physical investigations of the same observable}
\label{sec:physical}

The branches offer different ways to expose structure in the same
source and signed integral. A positive state space, probability law,
or channel is useful only if its controlled quantity can be connected
to the original determinant or to \eqref{eq:curvature}.
Table~\ref{tab:branches} gives the conceptual division of work. The
numbering agrees with the investigation folders.

\begin{table}[htbp]
\centering\small
\renewcommand{\arraystretch}{1.18}
\begin{tabularx}{\textwidth}{@{}p{.07\textwidth}>{\raggedright\arraybackslash}p{.28\textwidth}X@{}}
\toprule
Branch & Viewpoint & Quantity or structure to be understood\\
\midrule
01 & Circle thermal trace & Complete source and its spectral derivatives.\\
02 & Fermionic correlators & Signed overlaps and all finite occupation amplitudes.\\
03 & Connected Mellin response & Complete logarithmic curvature and saddle localization.\\
04 & Biorthogonal ensembles & Positive carrier with its signed insertion retained.\\
05 & Schur and Toda structure & Coefficient identities, evolution, and oscillator geometry.\\
06 & Topological field theory & Whether the required insertion is protected.\\
07 & Donaldson--Witten / Vafa--Witten & A specific four-dimensional source dictionary and boundary response.\\
08 & Open quantum dynamics & Conditional evolution realizing the same signed amplitudes.\\
\bottomrule
\end{tabularx}
\caption{The shared target is a signed comparison; the branches organize different mathematical interfaces to it.}
\label{tab:branches}
\end{table}

\subsection{01: a circle thermal trace supplies the exact source}

A particle on a circle of circumference $2\pi$, with Hamiltonian
$L=-\pi\,d^2/dx^2$, has energies $\pi m^2$. Its thermal trace is
$\Tr e^{-\beta L}=\vartheta(\beta)$. Combining this with
\eqref{eq:primitive} gives
\begin{equation}
 \Phi(u)=\beta^{1/4}\Tr\left[
   \bigl(\beta^2L^2-\tfrac32\beta L\bigr)e^{-\beta L}\right],
 \qquad \beta=e^{4u}. \label{eq:thermal}
\end{equation}
For $u\ge0$, the spectral filter is positive on every nonzero mode;
the zero mode is annihilated. This reconstructs the complete source,
including both terms of \eqref{eq:source}, without an arithmetic
truncation.

The thermal description suggests energy cumulants and source identities.
Its limitation is equally precise: differentiating the source and
forming $\Delta_r$ does not turn the derivative orbitals into identical
vectors. The accepted branch assessment retains a signed cumulant
residual. Positivity of the thermal trace alone does not control the
Mellin response \evid{P1}.

\subsection{02: fermions express determinants as amplitudes}

On $L^2((0,\infty),du)$ take the real orbitals
\begin{equation}
 f_i(u)=\frac{\Phi^{(2i)}(u)}{\sqrt{\Phi(u)}},\qquad
 g_{k,j}(u)=\frac{u^{2(k+j)}\sqrt{\Phi(u)}}{(2(k+j))!}.
\end{equation}
Their pairing is $\ip{f_i}{g_{k,j}}=c_{k+j-i}$ by \eqref{eq:ibp}.
The canonical anticommutation relations therefore give
\begin{equation}
 D_{r,k}=\bra{0}a(f_{r-1})\cdots a(f_0)
                  a^\dagger(g_{k,0})\cdots a^\dagger(g_{k,r-1})\ket{0}.
 \label{eq:fermion}
\end{equation}
Equivalently this is the inner product of two Slater wedges.
The annihilator order fixes the determinant sign. The elementary
fermion identities and Schur-function connections are discussed in
\cite{AlexandrovZabrodin2013}; the theta orbitals are the local
construction \evid{M3,P2}.

The two wedges differ, so the amplitude is not a squared norm. A
positive Hilbert-space metric yields a Cauchy--Schwarz magnitude bound,
not its sign. With $b_j=u^{2j}\sqrt\Phi$, let $\ket{F_r}=f_0\wedge\cdots\wedge f_{r-1}$
and $\ket{B_r}=b_0\wedge\cdots\wedge b_{r-1}$. The same construction gives
$M_r(s)=\bra{F_r}e^{s\widehat T}\ket{B_r}$, where
$\widehat T=2\sum\log u_j$ acts by multiplication. The logarithmic
derivatives are connected \emph{cross} matrix elements, not ordinary
positive-state variances.

There is also one bounded lattice model encoding every finite minor.
On $\ell^2(\N_0)$ let $S\ket0=0$ and $S\ket j=\ket{j-1}$. For $v>0$,
\begin{equation}
 K_v=\frac{F(vS)}{F(v)},\qquad \|K_v\|\le1,
 \quad
 \bra I\Gamma(K_v)\ket J=
 \frac{v^{\sum J-\sum I}}{F(v)^r}\det[c_{j_b-i_a}].
 \label{eq:transfer}
\end{equation}
Here $I,J$ are increasing occupation lists of the same length $r$,
and $\Gamma(K_v)$ is the direct sum of its exterior powers, with
vacuum amplitude $+1$. The norm bound follows from $c_n\ge0$ and
$\|S\|=1$. Nonnegativity of all these ordered amplitudes is exactly
the coefficient condition in \eqref{eq:rh}, not a weaker replacement.

On $\operatorname{span}\{\ket0,\ket1\}$ the transfer has a nonzero
Jordan off-diagonal $vc_1/F(v)$ and repeated diagonal $c_0/F(v)$.
It cannot be made self-adjoint by a positive definite inner product on
that invariant block. This rules out that particular Hermitian
interpretation of $K_v$; it does not rule out total positivity or a
larger Hermitian environment \evid{P2}.

\subsection{03: the curvature is a connected response}

The Mellin branch treats $s$ as a source for $T$ and asks directly for
the connected response \eqref{eq:curvature}. It is the analytical
consumer of the other descriptions. Its completed local saddle
calculation resolves cancellations in the derivative determinant and
identifies a three-particle Gaussian limit. Section~\ref{sec:saddle}
states the resulting full-integral obligation. A local model by itself
does not establish a curvature bound on any complete tail \evid{P3}.

\subsection{04: a positive ensemble with a signed insertion}

The carrier in \eqref{eq:carrier} is a biorthogonal ensemble. In the
full orthant set
\[
 a_i(u)=u^{2i+s}\sqrt{\Phi(u)},\qquad
 b_i(u)=R(u)^i u^s\sqrt{\Phi(u)},\qquad 0\le i<r.
\]
The product of the two orbital determinants is $w_{r,s}$, and its
normalized full-orthant density is $w_{r,s}/(r!Z_r(s))$.
If $H_{ij}=\int a_i b_j$, then $\det H=Z_r(s)>0$.
The correlation kernel is
\[
 \mathsf K(x,y)=\sum_{i,j}a_i(x)(H^{-1})_{ji}b_j(y).
\]
Its determinantal particle correlations describe a positive ensemble
\cite{Borodin1998,DesrosiersForrester2006}. The target remains
$M_r=Z_r\E h_r$: it has an insertion that the ensemble's positivity
does not sign.

Source integration by parts suggests Ward identities, meaning exact
relations among responses induced by a change of integration variables.
The worked dilation identity in \evid{P4} retains derivatives of the
insertion and leaves a mixed response unbounded. The selected finite
polynomial closure fails in its stated class. This is useful information
about that closure, not an exclusion of every ensemble estimate.

\subsection{05: Schur algebra, Toda identities, and a bosonic state}

Define formal times by
\begin{equation}
 \frac{F(w)}{c_0}=\exp\left(\sum_{m\ge1}t_mw^m\right).
\end{equation}
This is a formal logarithm near zero, not a global zero-free logarithm.
Jacobi--Trudi gives
\begin{equation}
 D_{r,k}=c_0^r s_{(k^r)}(\mathbf t), \label{eq:schur}
\end{equation}
where $(k^r)$ is a rectangle with $r$ rows and $s_\lambda$ is the
Schur function at this specialization. Desnanot--Jacobi is the
associated bilinear determinant relation. Schur positivity at an
arbitrary formal specialization is not automatic; representation
dimensions cannot be substituted for \eqref{eq:schur}.

Source heat, $\Phi_\tau=e^{\tau u^2}\Phi$, has exact coefficient flow
\begin{equation}
 \partial_\tau c_n(\tau)=(2n+2)(2n+1)c_{n+1}(\tau).
 \label{eq:sourceheat}
\end{equation}
The branch's finite-prefix nonclosure result says that the tested
rectangular-ratio flow cannot be determined from just a fixed finite
coefficient prefix over its stated larger class of sources. It does
not exclude an identity using the full theta source \evid{P5}.

A second construction is the normalizable one-oscillator state
\begin{equation}
 \ket{\psi_\theta}=\mathcal N^{-1}
       \sum_{n\ge0}\sqrt{(2n)!}\,c_n\ket{2n},\qquad
 \mathcal N^2=\int_0^\infty\!\int_0^\infty
                 \Phi(u)\Phi(v)\cosh(uv)\dd u\dd v<\infty.
\end{equation}
Its Bargmann function is $F(z^2)/\mathcal N$; hence the zero-location
question in \eqref{eq:rh} becomes confinement of its stellar zeros to
the imaginary axis. The associated Husimi probability is nonnegative
for every state; its zero confinement is the substantive condition.
The stellar terminology follows \cite{ChabaudMarkhamGrosshans2020}.
For the unnormalized state
$\chi_\tau=\sum_{n\ge0}\sqrt{(2n)!}c_n(\tau)\ket{2n}$,
\eqref{eq:sourceheat} gives $\partial_\tau\chi_\tau=a^2\chi_\tau$,
where $a$ is the bosonic annihilation operator. This equality holds on
the actual analytic-vector orbit for each finite real $\tau$. It is not unitary squeezing:
the adjoint creation term is absent. The proposed attenuated
implementation remains a separate physical question \evid{M2,P5}.

\subsection{06: protection must apply to the actual insertion}

Supersymmetric or topological language is valuable when a symmetry
protects the observable being studied. In circle de Rham quantum
mechanics, take $H=\pi\Delta$, $Q=\sqrt\pi\,d$, and form degree $N$.
The excited degree-fugacity supertrace is
\begin{equation}
 \Tr\bigl((-1)^N y^N P_{\rm exc}e^{-\beta H}\bigr)
        =(1-y)(\vartheta(\beta)-1).
\end{equation}
At $y=1$ the protected excited supertrace vanishes. Recovering theta
requires its response, with insertion $\mathcal O=-NP_{\rm exc}$.
It obeys $[Q,\mathcal O]=Q$, so that insertion is not protected.
The same operation that would erase the unprotected response would
erase the source needed for the coefficient problem \evid{P6}.
This is a concrete supersymmetric model and a scoped failure of a
protection argument; it is not a Donaldson-invariant identification.

\subsection{07: four-dimensional theories and a completed anomaly}

The Donaldson--Witten investigation first tested a specified
$\mathcal N=2$ dictionary on $\mathbb{CP}^2$. Selection rules and
coefficient relations obstruct the tested protected-correlator match,
including the extension by identity plus a finite odd polynomial of
point insertions. This statement fixes the manifold, bundle, insertion
family, and rescaling dictionary. It is not a no-go theorem for
Donaldson theory \evid{G1}.

A different construction uses the completed $\mathcal N=4$,
$SO(3)$ Vafa--Witten partition function in the $w_2=0$ sector.
The anomaly formula of Dabholkar--Putrov--Witten
\cite[eq.~(2)]{DabholkarPutrovWitten2020} is
\begin{equation}
 \partial_{\bar\tau}Z_0
 =\frac{3}{16\pi i\,y^{3/2}\eta(\tau)^3}
                 \sum_{m\in\mathbb Z}\bar q^{m^2},
 \qquad y=\operatorname{Im}\tau,\quad q=e^{2\pi i\tau}.
\end{equation}
Take the antiholomorphic derivative before restricting to
$\tau=i\beta/2$. The normalized response
\begin{equation}
 \mathcal A(\beta)=\frac{16\pi i}{3}(\beta/2)^{3/2}
 \eta(i\beta/2)^3
 \left.\partial_{\bar\tau}Z_0\right|_{\tau=i\beta/2}
 =\vartheta(\beta) \label{eq:anomaly}
\end{equation}
recovers the complete integer-square sum. Physical normalization and
orientation qualifications remain those recorded in \evid{G1}; the
explicit mathematical completion fixes the convention here.

Its completed coupling Mellin period is
\begin{equation}
 \mathcal X_{\rm VW}(z)=\frac12+\frac{z(z-1)}4
 \int_1^\infty[\mathcal A(\beta)-1]
       [\beta^{z/2}+\beta^{(1-z)/2}]\frac{d\beta}{\beta}
 =\xi(z). \label{eq:period}
\end{equation}
The equality follows by splitting the theta Mellin integral at one
and using theta inversion. It is a nonlocal transform over coupling,
not an equality with a local protected correlator. The zero-mode
subtraction and additive $1/2$ preserve the values at $z=0,1$.
The anomaly derivative removes holomorphic contributions, so their
counting interpretation does not by itself transmit an inequality to
the extracted response.

The subsequent bounded Ward analysis expresses the same signed
determinant through differentiated boundary fluxes and leaves
\begin{equation}
 MM''-(M')^2-\kappa_3M^2
\end{equation}
uncontrolled. It gains no comparison interval. That particular route
is parked; other identities or dictionaries remain outside the
assessment \evid{G2}.

\subsection{08: conditional open-system dynamics}

The transfer $K_v$ in \eqref{eq:transfer} can be studied as a
conditional no-jump evolution. Where it is invertible, its generator is
\begin{align}
 G(v)&=K_v'K_v^{-1}
      =SF'(vS)F(vS)^{-1}-\frac{F'(v)}{F(v)}I,\notag\\
 R(v)&=-\bigl(G(v)+G(v)^\dagger\bigr),\qquad
 H(v)=\tfrac i2\bigl(G(v)-G(v)^\dagger\bigr).
 \label{eq:nj}
\end{align}
If $R(v)\ge0$, then $G=-iH-R/2$ admits a loss realization.
On fermionic sectors, a quadratic Hamiltonian $d\Gamma(H)$ and linear
loss operators with summed quadratic rate $d\Gamma(R)$ give the
no-count propagator $\bigwedge^r K_v$. Averaging all jump records
produces density-matrix evolution; it is a different object from the
selected conditional amplitude. This distinction is standard in the
quantum-trajectory description \cite{CastinDalibardMolmer08054002v1}.

The latest completed MC08 construction proves admissibility on
$0\le v\le1/25$, a finite-mode count instrument with an explicit
collision-limit error, a coherent readout retaining the amplitude's
sign, and a compatible vacuum-Fock extension. Its separate assessment
is same-model internal LLM review with shared premises \evid{O1}.
The microscopic limit is taken at fixed mode cutoff with fresh
environment modes. It is not an exact finite-resource laboratory
implementation, a bounded full-Fock generator theorem, or a new
determinant inequality.

An additional actual-theta argument in \evid{O2} rules out preservation
of the fixed nonnegative ordered occupation cone on \emph{every
intermediate activity interval}, already in the three-particle sector.
Its status is proof-author checking, with separate verification pending
in the inspected record. It does not produce a negative theta endpoint
minor or close the open-system program. Moreover,
\begin{equation}
 A_{3,k}(v)=\bra{0,1,2}\Gamma(K_v)\ket{k,k+1,k+2}
          =\frac{v^{3k}}{F(v)^3}D_{3,k}
\end{equation}
implies the exact cancellation
\begin{equation}
 \frac{A_{3,k-1}(v)A_{3,k+1}(v)}{A_{3,k}(v)^2}
 =\frac{D_{3,k-1}D_{3,k+1}}{D_{3,k}^2}. \label{eq:activity}
\end{equation}
Changing activity alone does not change this ratio. A useful dynamics
argument must add a source-dependent inequality, not infer it from
complete positivity or conditional normalization.

\section{The most explicit remaining analytical task: localization}
\label{sec:saddle}

The connected-response calculation \evid{P3} offers a concrete
large-$s$ route. It also shows the distinction between a plausible
local model and a global estimate. Let
\begin{equation}
 s=2\pi a e^{4a},\qquad t=\pi e^{4a},\qquad
 \sigma^2=\frac{a^2}{2s(1+4a)},\qquad u_j=a+\sigma x_j,
 \label{eq:scales}
\end{equation}
with $a>0$. The exponent $E_s(u)=2s\log u-\pi e^{4u}$ has its
stationary point at $a$, with $-E_s''(a)=\sigma^{-2}$.

For the complete signed density $f_s$ in \eqref{eq:twocopycurv}, set
\begin{align}
 K(s)&=(2e^at^2)^3e^{3E_s(a)}(256at^2)^3\sigma^9,\notag\\
 \rho_s(\mathbf x)&=K(s)^{-1}\sigma^3
             f_s(a+\sigma x_1,a+\sigma x_2,a+\sigma x_3),\notag\\
 \mathcal D_s&=\{-a/\sigma<x_1<x_2<x_3\}. \label{eq:rhos}
\end{align}
On every fixed ordered box, the accepted local limit is
\begin{equation}
 \rho_s(\mathbf x)\longrightarrow
 \rho_G(\mathbf x)=e^{-\frac12\sum x_j^2}V(\mathbf x)^2.
 \label{eq:local}
\end{equation}
The factor $\sigma^9$ includes three powers from volume and six from
the two Vandermondes. It matters: replacing every moment entry by its
leading scalar asymptotic can leave a rank-one matrix and erase the
leading determinant. The divided-difference calculation resolves this
cancellation before estimating the integral.

Let $G_0=2(2\pi)^{3/2}$ be the ordered Gaussian integral, and define
\begin{equation}
 Y_s=\frac a\sigma\sum_j\log(1+\sigma x_j/a),\qquad
 I_p(s)=G_0^{-1}\int_{\mathcal D_s}Y_s^p\rho_s\dd\mathbf x,
 \quad p=0,1,2.
\end{equation}
These are complete \emph{signed} integrals, with $I_0=M/(G_0K)>0$.
Centering $T$ algebraically at each fixed $s$ gives the exact identity
\begin{equation}
 \mathcal C_3(s)=\frac{2}{s(1+4a)}
       \left(\frac{I_2}{I_0}-\left(\frac{I_1}{I_0}\right)^2\right).
 \label{eq:susceptibility}
\end{equation}
This does not differentiate an uncontrolled asymptotic remainder.

\begin{lemma}[A sufficient tail payment]
If, for every $s\ge S\ge4\pi e^8$, one proves
\begin{equation}
 I_0(s)\ge\tfrac12,\qquad I_2(s)\le6, \label{eq:payments}
\end{equation}
then $\mathcal C_3(s)\le\kappa_3(s)$ on that tail.
\end{lemma}
\begin{proof}
The nonnegative square in \eqref{eq:susceptibility} gives
$\mathcal C_3\le24/[s(1+4a)]$. The threshold implies $a\ge2$ and
$s\ge20$, while
$\kappa_3(s)\ge\sum_{j=0}^2(s+j+1/2)^{-1}\ge3/(s+5/2)$.
Finally $24/(9s)\le3/(s+5/2)$ for $s\ge20$.
\end{proof}

The lemma proves an implication, not its theta hypotheses. In the
fixed box $\mathcal B_{12}=\{-12\le x_1\le x_2\le x_3\le12\}$,
write $I_p=A_p+B_p$, where $B_p$ is the full complementary integral.
One sufficient set of bounds is
\begin{equation}
 A_0\ge\tfrac34,\qquad A_2\le4,\qquad
 B_0\ge-\tfrac14,\qquad B_2\le2. \label{eq:fourpayments}
\end{equation}
The first two follow eventually from the fixed-box limit, without an
effective threshold in the existing result. The last two remain
unproved. They include configurations with just one or two nodes far
from the saddle, the negative origin-side region, and all spatial
tails. A bound on the source along a common ray does not cover these
relative integrals.

The Gaussian source of \eqref{eq:gaussM} and the Gaussian node density
of \eqref{eq:local} are different comparisons. The former gives an
exact all-$s$ curvature; the latter is a local asymptotic of theta.
Neither implies full-domain convergence in the other setting.
Even a proof of \eqref{eq:payments} on an explicit tail would leave
$[0,S]$, especially $[0,2]$, to be controlled.

\section{What the subinvestigations must preserve and establish}
\label{sec:obligations}

\subsection{Parameters and observables}

\begin{table}[htbp]
\centering
\renewcommand{\arraystretch}{1.15}
\begin{tabularx}{\textwidth}{@{}p{.16\textwidth}X@{}}
\toprule
Parameter & Meaning and required distinction\\
\midrule
$r,k$ & Determinant size and coefficient index; arbitrary-order induction requires both quantifiers.\\
$d,n$ & Jensen degree and shift; they are not $r,k$.\\
$s$ & Real Mellin tilt of the complete determinant integral.\\
$u,\beta$ & Source coordinate and inverse temperature, related by $\beta=e^{4u}$.\\
$\tau$ & Source-heat parameter in \eqref{eq:sourceheat}; in Section 5.7 only, the conventional complex gauge coupling.\\
$v$ & Transfer activity in \eqref{eq:transfer}; not a Mellin or source-heat parameter.\\
\bottomrule
\end{tabularx}
\caption{Parameter changes need an explicit bridge to the target observable.}
\end{table}

Each new estimate should identify its exact source, normalization,
domain, and lower-order positivity assumptions. In particular, replacing
an amplitude by its squared modulus, replacing $h$ by $|h|$, or
retaining only a positive carrier changes the target. Such quantities
may still be useful, but they need a proved comparison back to the
original signed expression.

There are two legitimate kinds of progress. A mathematical advance
proves a new bound on a stated domain, resolves a limiting or boundary
obligation, or excludes a precisely defined mechanism. A physical
advance gives a useful exact model, instrument, or observable dictionary.
The latter need not yet improve positivity, but its value should be
stated in those terms. The eight branches have achieved several
results of the second kind and several scoped obstructions; the
global Gaussian comparison remains open.

\subsection{Controls test premises, not labels}

The repository control inventory and methodological index
\evid{C1,C2} impose a claim-specific screen. Applied to the present
program, it has the following content.
\begin{itemize}
\item The Gaussian source checks normalization and the equality case.
It has $h_r=1$ and is not the canonical theta source.
\item The two-point relaxation retains association and increasing
averages but fails the canonical derivative-determinant realization.
It disproves an association-only deduction, not the theta hypothesis.
\item Euler-factor controls preserve positivity and evenness but change
arithmetic reconstruction. Their curvature comparison is untested here;
a positive-time collision alone does not refute this time-zero pilot.
\item Weighted-square and finite-prime perturbations alter the canonical
weights. In the relevant finite-endpoint constructions they also fail
full reflection, so finitely many matched endpoints cannot justify
\eqref{eq:ibp} at every order.
\item Reflected super-Gaussian and smooth positive-primitive controls
can preserve substantial shape information while failing literal theta
reconstruction. Generic determinant and response identities may survive;
no uncomputed sign is assigned to their quantitative comparison.
\item The polynomial $F_*(w)=1+w+2w^2$ has positive coefficients and
a contraction transfer, but $c_1^2-c_0c_2=-1$. MC08's checked local
completion retains this negative two-fermion amplitude. Generic channel
positivity therefore cannot supply ordered amplitude positivity.
\end{itemize}
Merely naming the complete theta source excludes these modified sources
at the level of identity. It is not yet a proof that a proposed
inequality uses enough of that identity to distinguish them.

\subsection{The order-four pilot is not yet an induction theorem}

A successful pilot would prove \eqref{eq:pilot} for all $k\ge1$,
with the anchor handled separately. To iterate, the program needs a
corresponding estimate for every arbitrary finite $r$. Constants may
depend on $r$; a finite list of orders is insufficient. If the continuous
route is used, one additionally needs $M_r(s)>0$ for all required real
tilts, not just integer $s$.

The accepted order-three monotonicity cannot simply be reused at order
four. The next insertion introduces $g_3=\Phi^{(6)}/\Phi$, and a
canonical separated-node obstruction to unchanged $h_4$ coordinatewise
monotonicity is recorded in the sibling signed-averaging assessment
\evid{J3}. That obstruction does not prove a negative integrated
average. A discrete proof of \eqref{eq:G} could avoid the extra
real-interpolation premise, and an integrated comparison could survive
a failure of the pointwise envelope.

The immediate analytical task is therefore well defined: obtain one
complete-source bound on the signed response, with every compact,
endpoint, mixed-node, and tail contribution accounted for. The source
identities and physical descriptions explain where such information
might come from. The missing estimate is the step that would turn those
representations into an order-raising argument.

\appendix
\section{Repository source guide and evidence status}
\label{sec:sources}

This manuscript synthesizes the investigation README, its research
overview, and the records below. Paths in the first list are relative to
\path{papers/investigations-5/07_positive_theta_primitives/01_jensen/01_mellin_curvature/}.
The accompanying integration note contains exact paths and input hashes.
The individual research-progress registers remain authoritative for task
status after this dated draft.

\begin{description}[style=nextline,leftmargin=2.5em]
\item[\evid{J1}] Jensen Note14, \emph{Higher order signed averaging and radial
flow feasibility}, 1 November 2026, in \path{../notes/}. It supplies the
all-order integral, endpoint conventions, strict-minor bridge, and
scope of possible continuations. Same-assistant analytical checking.
\item[\evid{J2}] Jensen Review01, \emph{Global sign certificate and
signed-averaging assessment}, in \path{../reviews/}. Accepts the complete
source signs and order-three normalized index comparison. Same-assistant
and same-model internal checking with inherited premises; no specialist review.
\item[\evid{J3}] Jensen Review02, \emph{Parallel branch evaluation review},
in \path{../reviews/}, together with the sibling signed-averaging
initial return. Records the scoped canonical $h_4$ monotonicity obstruction.
\item[\evid{M1}] \path{notes/01_INITIAL_BRANCH_EVALUATION_20261101.md}.
Exact factorial budget, real-tilt extension, Gaussian comparison,
two-point relaxation, and endpoint/tail requirements. Internally checked
in the parent comparative assessment.
\item[\evid{M2}] \path{notes/06_PARALLEL_PHYSICAL_DIRECTIONS_SYNTHESIS_20261101.md}
and \path{reviews/01_PARALLEL_PHYSICAL_DIRECTIONS_ASSESSMENT_20261101.md}.
Six completed physical assessments and their separately checked scopes.
Conceptual authors: gpt-6-astra/ultra; separate reviewers:
gpt-6.1-sol/ultra. These are internal LLM reviews, not specialist validation.
\item[\evid{M3}] \path{notes/04_FERMIONIC_CORRELATORS_AND_PHYSICAL_INTERPRETATIONS_20261101.md}.
The elementary orbital, ensemble, and Schur dictionaries. Same-assistant
checking; later branch results take precedence over its early status wording.
\end{description}

The branch identifiers below refer to their local Note01 and separate
Review01, unless a later record is specified. Each branch folder contains
the full bibliography routes, derivations, limits, and progress register.
\begin{description}[style=nextline,leftmargin=2.5em]
\item[\evid{P1}] \path{01_circle_thermal_trace/}: complete spectral filter and
the residual in its proposed source comparison.
\item[\evid{P2}] \path{02_fermionic_correlators/}: bounded all-sector transfer,
normalizations, and the fixed-positive-metric Jordan obstruction.
\item[\evid{P3}] \path{03_connected_mellin_response/}: complete signed saddle,
$\sigma^9$ normalization, exact susceptibility, and unpaid global budgets.
\item[\evid{P4}] \path{04_biorthogonal_ensembles/}: inserted Ward identity
and the selected finite-polynomial-closure obstruction.
\item[\evid{P5}] \path{05_schur_toda_structure/}: source-heat evolution,
finite-prefix nonclosure, and the normalizable oscillator state.
\item[\evid{P6}] \path{06_topological_field_theory/}: specified supersymmetric
circle theory and its unprotected source response.
\item[\evid{G1}] \path{07_donaldson_witten_theory/}, Notes01--02 and
Reviews01--02: scoped $\mathcal N=2$ coefficient obstructions and the
completed $\mathcal N=4$ anomaly dictionary. Separate cross-model
internal LLM checking, with physical-normalization qualifications.
\item[\evid{G2}] The same folder, Note03, \emph{Boundary Ward feasibility
return}: unchanged signed residual and no new comparison domain.
Proof-author checking only.
\item[\evid{O1}] \path{08_open_quantum_dynamics/}, Notes01--03 and Review02,
\emph{Open quantum dynamics assessment}: accepted local interval,
instrument, readout, and convergence at stated scope. Separate same-model
internal LLM assessment with shared inputs.
\item[\evid{O2}] The same folder, Note04, \emph{Signed comparison decision},
and Review03: scoped every-intermediate-interval cone obstruction and
activity cancellation. Proof-author checking only; separate E04
verification pending in the inspected register. Note05 preserves the
subsequent broader discussion without a new proved estimate.
\end{description}

Two repository-wide guides govern interpretation:
\begin{description}[style=nextline,leftmargin=2.5em]
\item[\evid{C1}] \path{result-summaries/README.md}, especially cluster 12,
\emph{Collision controls and screening}, and cluster 14,
\emph{Theta identities and arithmetic bookkeeping}. These route to the
actual proofs and certificates; their summaries do not certify results.
\item[\evid{C2}] \path{avoid/README.md}, especially patterns 01
(target reformulation), 02 (control premises), 07 (scoped exclusions),
12 (joint signed estimates), 13 (a factor is not the target), and 19
(status files that lag the findings).
\end{description}

\section*{Acknowledgements and preparation}
This manuscript was drafted for Edward Baker with substantial assistance
from OpenAI's Codex, using the recorded model \texttt{gpt-6-astra} at
reasoning effort \texttt{max}. The assistance included source synthesis,
exposition, algebraic consistency checks, and typesetting. The underlying
research has the distinct model and review provenance recorded in its
source files. The present manuscript received proof-author consistency
and layout checking, not independent mathematical or specialist review.
No historical novelty claim is made for the classical identities or the
physical formulations assembled here.

\bibliographystyle{plainurl}
\bibliography{manuscript_references}
\end{document}
